\documentclass[10pt,twocolumn]{article}

\usepackage[utf8]{inputenc}
\usepackage[T1]{fontenc}
\usepackage{lmodern}
\usepackage[margin=0.72in,columnsep=0.28in]{geometry}
\usepackage{microtype}
\usepackage{graphicx}
\usepackage{booktabs}
\usepackage{array}
\usepackage{amsmath,amssymb}
\usepackage{xcolor}
\usepackage{hyperref}
\usepackage{enumitem}
\usepackage{caption}
\usepackage{subcaption}
\usepackage{fancyhdr}
\usepackage{titlesec}
\usepackage{balance}

\definecolor{sqblue}{HTML}{0B6E8E}
\definecolor{sqgreen}{HTML}{087F5B}
\definecolor{sqred}{HTML}{C92A4A}
\definecolor{sqgray}{HTML}{4A5568}
\hypersetup{
  colorlinks=true,
  linkcolor=sqblue,
  citecolor=sqgreen,
  urlcolor=sqblue,
  pdftitle={SpectraQuant: From Budget-Aware Allocation Failure to Verified Clipping Transfer},
  pdfauthor={Oussema Harrabi}
}
\graphicspath{{../../reports/figures/}{../figures/}}
\captionsetup{font=small,labelfont=bf}
\setlist{nosep,leftmargin=*}
\setlength{\parindent}{0.85em}
\setlength{\parskip}{0.1em}
\titleformat{\section}{\large\bfseries\color{sqblue}}{\thesection}{0.55em}{}
\titleformat{\subsection}{\normalsize\bfseries\color{sqgray}}{\thesubsection}{0.5em}{}
\pagestyle{fancy}
\fancyhf{}
\fancyhead[L]{\small SpectraQuant Technical Report}
\fancyhead[R]{\small October 2026}
\fancyfoot[C]{\thepage}
\renewcommand{\headrulewidth}{0.3pt}

\newcommand{\ppl}{\mathrm{PPL}}
\newcommand{\bytes}{\mathrm{bytes}}
\newcommand{\good}[1]{\textcolor{sqgreen}{\textbf{#1}}}
\newcommand{\bad}[1]{\textcolor{sqred}{\textbf{#1}}}
\newcommand{\sqfigure}[2]{%
  \IfFileExists{#1}{\includegraphics[width=#2]{#1}}{%
    \fbox{\parbox[c][0.16\textheight][c]{0.94#2}{\centering
      Figure asset: \texttt{\detokenize{#1}}}}}}

\title{\vspace{-1.2em}\textbf{SpectraQuant: From Budget-Aware Allocation Failure\\
to Verified Clipping Transfer}\\[0.35em]
\large A Reproducible Study of Local Compression Objectives and Whole-Model Quality}
\author{Oussema Harrabi\\
\small Independent Research Project\\
\small \url{https://github.com/OussemaHarrabi/SpectraQuant}}
\date{12 October 2026}

\begin{document}
\maketitle

\begin{abstract}
SpectraQuant investigates whether inexpensive local compression objectives can preserve
autoregressive language-model quality under a measured equal-byte checkpoint budget.
The project began with a sensitivity-aware joint low-rank and low-bit allocator.
On SmolLM2-135M, its trained arm failed to beat uniform INT4: perplexity was
23.1698 without the preparation regularizer and 23.3131 with it, versus 18.6193
for uniform INT4 at matched target bytes. A deployment-matched follow-up then
evaluated the serialized and reloaded residual representation; validation again
retained uniform INT4, while six alternatives worsened held-out perplexity by
5.52--14.75\%. Diagnostic experiments showed that summed marginal layer gains
overpredicted the quality of the jointly modified model.

This failure motivated a simpler conventional clipping control. A first frozen
C4 holdout improved from 28.0056 to 27.4656 perplexity (1.9279\%) at identical
payload and full-checkpoint bytes. A subsequently frozen transfer study covered
SmolLM2-360M and Pythia-410M on reserved C4 and PG19 prefixes. The primary
weight-MSE control improved 2/4 cells and achieved a 0.5752\% geometric-mean
reduction, but failed its predeclared 3/4 consistency gate. Predeclared
activation-weighted axis1 controls reduced C4 perplexity by 7.39--7.70\% on
SmolLM2 and 5.75--6.09\% on Pythia at identical bytes, while regressing on
Pythia PG19. The study therefore establishes a bounded C4 result and a
model/domain transfer failure, not a universal compression improvement.
All 52 expected Phase-4 artifacts and 42 checkpoint files were authenticated.
\end{abstract}

\noindent\textbf{Keywords:} transformer compression; post-training quantization;
INT4; clipping; sensitivity; low-rank approximation; reproducibility; negative results

\section{Introduction}

Transformer compression is often optimized through local quantities: weight
reconstruction error, activation-weighted error, Hessian surrogates, spectral
statistics, or layer-wise calibration loss. The operational outcome, however,
is global. Every modified layer participates in one autoregressive computation,
and small locally favorable changes may interact poorly when composed.

SpectraQuant asks one practical question:
\begin{quote}
\emph{Can cheap local compression scores choose a low-bit representation that
preserves end-to-end language quality at the same measured checkpoint budget?}
\end{quote}

The project is unusual in treating a failed original method as the beginning of
the investigation rather than an outcome to omit. The experimental progression
is:
\begin{enumerate}
  \item test a joint low-rank/low-bit allocator against uniform INT4;
  \item correct the evaluation to score the actual packed/reloaded operator;
  \item test whether marginal scores predict joint-model loss;
  \item reduce the intervention to conventional clipping;
  \item freeze a two-model/two-domain transfer study before reading holdouts.
\end{enumerate}

The final contribution is therefore methodological and empirical. The repository
implements deterministic experiment plans, equal-byte gates, packed-storage
checks, cloud notebook generation, authenticated collection, generated figures,
and explicit measurement classes. Empirically, it documents where local
objectives fail, where clipping helps, and where that help stops transferring.

\begin{figure*}[t]
  \centering
  \sqfigure{spectraquant-research-storyline.png}{\textwidth}
  \caption{The coherent study progression. Later phases simplify or test an
  assumption exposed by the preceding failure; they are not an unrelated search
  for a favorable number.}
  \label{fig:storyline}
\end{figure*}

\section{Background and Related Work}

Post-training quantization (PTQ) maps dense weights to a small discrete codebook
after training. Round-to-nearest (RTN) is a strong transparent baseline.
Second-order and activation-aware methods such as GPTQ~\cite{gptq},
AWQ~\cite{awq}, and SmoothQuant~\cite{smoothquant} demonstrate that calibration
statistics can materially change low-bit behavior. HAWQ-style methods use
sensitivity estimates to allocate precision~\cite{hawq}.

Low-rank adaptation and quantization are also frequently combined. LoftQ
initializes low-rank adapters against quantization error~\cite{loftq}, while
LR-QAT embeds low-rank structure inside quantization-aware training~\cite{lrqat}.
These works motivate the original SpectraQuant hypothesis but do not remove the
need for equal-storage comparisons or deployment-matched evaluation.

The present study does not claim a new PTQ algorithm. Its strongest evidence is
about \emph{evaluation validity}: local proxy improvements and lower weight
error can fail to predict whole-model perplexity, and a control that helps one
architecture/domain pair can regress on another.

\section{Problem Formulation}

For compressible layer $\ell$ with weight
$W_\ell\in\mathbb{R}^{d_{\mathrm{out}}\times d_{\mathrm{in}}}$, the original
allocator chooses rank $r_\ell$ and bit width $b_\ell$:
\begin{equation}
\min_{\{r_\ell,b_\ell\}}\sum_{\ell=1}^{L}\widehat{E}_{\ell}(r_\ell,b_\ell)
\quad\mathrm{s.t.}\quad
\sum_{\ell=1}^{L}\bytes_\ell(r_\ell,b_\ell)\leq B .
\label{eq:allocation}
\end{equation}
Here $\widehat{E}_\ell$ is a local sensitivity surrogate and $B$ is a frozen
payload budget. The deployed residual form is
\begin{equation}
\widehat{W}=Q(W)+Q(B)Q(A),
\end{equation}
with rank zero representing plain quantization. SpectraQuant measures the
serialized target payload and full checkpoint rather than inferring equality
from nominal bit widths.

The later clipping study fixes INT4/group32 quantization and chooses a clipping
scale $\alpha$ from a predeclared grid. Weight-only MSE minimizes local
reconstruction error. The activation-weighted diagonal proxy is proportional to
\begin{equation}
\widehat{E}_{\mathrm{diag}} =
\sum_j \mathbb{E}[x_j^2]\sum_i
\left(W_{ij}-Q_\alpha(W)_{ij}\right)^2 .
\label{eq:diagproxy}
\end{equation}
It omits cross-feature covariance and is not measured output error. To avoid a
degenerate comparison, activation weighting is evaluated on axis1 against
axis1 RTN and MSE controls; for axis0, the input second moment is constant
inside each optimized output-row group.

\section{Research Design}

\subsection{Decision discipline}

Three rules govern all headline comparisons:
\begin{enumerate}
  \item quality is compared only when measured storage matches;
  \item evaluation reloads the files that the experiment actually writes;
  \item model revisions, corpora, arms, gates and token budgets are frozen before
  fresh holdout access.
\end{enumerate}

Protocol changes are append-only amendments. Failed gates remain visible.
Exploratory calibration subsets are never relabeled as independent training
seeds, and exposed holdouts are not reused for tuning.

\subsection{Measurement classes}

The repository separates four forms of evidence:
\begin{itemize}
  \item \textbf{Class 1}: analytical size estimates;
  \item \textbf{Class 2}: dense execution of reconstructed quantized weights;
  \item \textbf{Class 3}: bytes of serialized payload/checkpoint files;
  \item \textbf{Class 4-CPU}: a bounded real CPU low-bit-kernel path.
\end{itemize}

Phase-4 perplexity is Class 2 and storage is Class 3. No Phase-4 number implies
GPU low-bit acceleration, service throughput, or latency.

\subsection{Phase-4 frozen transfer protocol}

The transfer design uses SmolLM2-360M and Pythia-410M at immutable revisions.
The reserved corpora are C4 validation rows 320--639 and the first 50 PG19 test
books as source pools. Calibration uses the first 200 C4 training documents.

Seven symmetric INT4/group32 checkpoints are written per model: RTN and
weight-MSE clipping on axis0 and axis1, plus activation-weighted axis1 clipping
from calibration subsets 0, 1, and 2. The clipping grid is
$\{1.00,0.95,0.90,0.85,0.80,0.75\}$. There is no holdout NLL search and no
gradient training. All checkpoints are frozen before holdout loading.

Each corpus/model cell predicts exactly 51,200 tokens with sequence length 1024
and preserved document boundaries. Comparisons are paired inside a model because
the two tokenizers differ. Document-cluster bootstrap intervals describe corpus
uncertainty; they are not training-seed intervals.

\begin{figure*}[t]
  \centering
  \sqfigure{spectraquant-evidence-architecture.png}{\textwidth}
  \caption{Evidence architecture. Frozen design and typed run specifications feed
  cloud execution; packed checkpoints are reloaded for evaluation; manifests,
  hashes and CI gates separate measured evidence from publication claims.}
  \label{fig:architecture}
\end{figure*}

\section{Experimental Progression}

\subsection{Original allocation result}

The equal-memory SmolLM2-135M comparison refuted the original improvement claim.
Uniform INT4 achieved 18.6193 perplexity. The trained allocated arm scored
23.1698 without the preparation regularizer and 23.3131 with it. The regularizer
therefore worsened the trained arm by approximately 0.62\%.

A bounded LR-QAT-direction check improved perplexity from 26.9432 to 24.1265.
This closed 21.8\% of the PTQ-to-FP16 gap, below the preregistered 25\% threshold.
It is reported as a partial directional reproduction rather than confirmation.

\subsection{Deployment-matched correction}

Phase 2 scored the serialized and reloaded $Q(W)+Q(B)Q(A)$ operator under an
equal INT4 target-payload ceiling. Validation retained uniform INT4. Six
mixed-bit or residual-inclusive alternatives worsened the held-out WikiText
perplexity by 5.52--14.75\%. Twenty-one checkpoint files passed validation.

\subsection{Joint-loss diagnosis and pivot}

Phase 3 compared summed marginal predictions with composed-model loss. All six
candidate allocations overstated joint calibration gains and worsened loss on
16 disjoint training blocks. This ruled out a simple explanation based only on
the earlier test split.

The intervention was then simplified to conventional weight-only clipping.
On a new C4 holdout, the fixed control improved perplexity from 28.0056 to
27.4656, a 1.9279\% reduction, with identical 59,719,680-byte target payloads
and 116,522,698-byte full checkpoints. This single-model pilot motivated the
frozen transfer test; it was not treated as proof of generality.

\begin{figure}[t]
  \centering
  \sqfigure{fixed-clipping-first-holdout.png}{\columnwidth}
  \caption{The first fresh-holdout clipping comparison. This bounded positive
  result triggered a transfer study rather than immediate generalization.}
  \label{fig:firstholdout}
\end{figure}

\section{Transfer Results}

\subsection{Primary weight-MSE result}

Table~\ref{tab:primary} compares MSE clipping with RTN using the same
group32/axis0 representation. Positive reduction denotes lower perplexity:
\begin{equation}
\Delta_{\ppl}=100\left(1-\frac{\ppl_{\mathrm{method}}}
{\ppl_{\mathrm{RTN}}}\right).
\end{equation}

\begin{table}[t]
\centering
\caption{Primary frozen transfer result.}
\label{tab:primary}
\small
\begin{tabular}{llrrr}
\toprule
Model & Corpus & RTN & MSE & $\Delta_{\ppl}$ \\
\midrule
Smol360 & C4   & 21.8140 & 22.3901 & -2.6408\% \\
Smol360 & PG19 & 16.6734 & 16.8819 & -1.2501\% \\
Pythia410 & C4   & 32.3251 & 31.3777 & +2.9311\% \\
Pythia410 & PG19 & 27.9244 & 27.0500 & +3.1313\% \\
\bottomrule
\end{tabular}
\end{table}

The within-cell log-ratio geometric-mean reduction is \good{0.5752\%}, exceeding
the 0.5\% magnitude requirement. Only \bad{2/4 cells} improve, below the frozen
3/4 consistency requirement. The primary transfer hypothesis therefore fails.

\begin{figure}[t]
  \centering
  \sqfigure{clipping-transfer-primary.png}{\columnwidth}
  \caption{Primary transfer cells and descriptive document-cluster intervals.}
  \label{fig:primary}
\end{figure}

\subsection{Activation-weighted secondary control}

Table~\ref{tab:secondary} reports the range across all three calibration subsets.
The baseline is each model's matched axis1 RTN checkpoint.

\begin{table}[t]
\centering
\caption{Activation-weighted PPL reduction ranges.}
\label{tab:secondary}
\small
\begin{tabular}{llr}
\toprule
Model & Corpus & Reduction range \\
\midrule
Smol360 & C4   & +7.3882 to +7.7028\% \\
Smol360 & PG19 & +5.8906 to +6.1849\% \\
Pythia410 & C4   & +5.7462 to +6.0939\% \\
Pythia410 & PG19 & -2.5608 to -1.0864\% \\
\bottomrule
\end{tabular}
\end{table}

The control improves reserved C4 on both architectures and SmolLM2 PG19, but
regresses on Pythia PG19. Calibration-domain sensitivity is plausible, yet this
study does not isolate it causally because it does not add book-domain
calibration. The supported statement is therefore a C4 improvement at equal
bytes, accompanied by a documented transfer failure.

\begin{figure}[t]
  \centering
  \sqfigure{clipping-transfer-controls.png}{\columnwidth}
  \caption{All predeclared clipping controls. Secondary results do not override
  the failed primary gate.}
  \label{fig:controls}
\end{figure}

\subsection{Storage equality and artifact verification}

All seven arms for a model have identical storage:
\begin{table}[h]
\centering
\caption{Measured Phase-4 storage per arm.}
\label{tab:storage}
\small
\begin{tabular}{lrr}
\toprule
Model & Target payload & Full checkpoint \\
\midrule
Smol360 & 176,947,200 B & 271,561,031 B \\
Pythia410 & 198,844,416 B & 302,615,137 B \\
\bottomrule
\end{tabular}
\end{table}

The collection process verified 52 expected artifacts and 42 dynamically
selected checkpoint files, including checksums, source revision, model/data
fingerprints, token counts, and the ordering constraint that checkpoint
selection precedes fresh holdout access.

\section{Why Local Objectives Failed}

Weight-MSE clipping reduced aggregate weight squared error by 9.98\% on
SmolLM2 axis0 and 10.19\% on axis1, yet worsened perplexity on both corpora.
Pythia axis1 supplies another counterexample: 10.69\% lower weight error
accompanied 7.06\% worse PG19 perplexity.

\begin{figure}[t]
  \centering
  \sqfigure{clipping-objective-vs-quality.png}{\columnwidth}
  \caption{Lower local reconstruction error is not sufficient for lower
  autoregressive perplexity.}
  \label{fig:objective}
\end{figure}

The Phase-3 allocation diagnostics show the same issue from another angle.
Layer effects measured in isolation were not additive after simultaneous
deployment. Transformer residual paths, normalization, activation shifts and
error propagation couple layers. A useful proxy must therefore be validated
against composed-model outcomes rather than only local score correlation.

\begin{figure}[t]
  \centering
  \sqfigure{allocation-prediction-vs-joint.png}{\columnwidth}
  \caption{Marginal allocation predictions versus joint-model behavior.}
  \label{fig:joint}
\end{figure}

\section{Reproducibility and Compute}

The Phase-4 scientific source was frozen at commit
\texttt{71a1ec11cea6be90a6ab49356e81d980f4667784}. Typed run specifications record
model revisions, dataset windows, arms, token budgets, clipping grids and
acceptance gates. The generated Kaggle notebooks preserve those inputs.

Two private Kaggle runs completed on their first submission. Their runner
intervals correspond to approximately 1.511184 visible T4 device-hours combined,
including installation but excluding export/platform overhead. This value is
execution telemetry, not exact account billing.

The project records Python, PyTorch, Transformers, CUDA and hardware metadata.
Collection authenticates remote manifests before accepting derived tables.
Scientific figures are generated from committed artifacts with input/output
hashes. Local CI checks formatting, types, tests, documentation and security.

\section{Limitations}

\begin{itemize}
  \item Phase-4 calibration subsets are not independent model-training seeds.
  \item Only two sub-billion-parameter model families and two reserved corpora
  are tested.
  \item Perplexity is measured through dense reconstruction, not a deployed
  low-bit GPU kernel.
  \item The diagonal activation proxy omits feature covariance.
  \item Modern PTQ baselines such as full GPTQ/AWQ implementations are reviewed
  but not run as matched Phase-4 arms.
  \item Exposed holdouts cannot support further tuning; a new hypothesis
  requires a new protocol and untouched data.
\end{itemize}

\section{Conclusion}

SpectraQuant did not validate its original joint low-rank/low-bit allocator.
Correcting deployment evaluation and inspecting joint loss strengthened that
negative result. The subsequent pivot showed that conventional clipping can
deliver meaningful equal-byte gains, but also that transfer depends on model,
corpus and grouping. Activation-weighted clipping produced repeatable C4
improvements across both tested architectures and all three calibration subsets,
while failing on Pythia PG19.

The main result is therefore not a universal compression recipe. It is a
reproducible account of how to turn a failed hypothesis into useful engineering
knowledge: measure the deployed representation, freeze the next test, retain
negative cells, authenticate artifacts, and bound every claim by the evidence
that actually ran.

\section*{Artifact Availability}

Code, protocols, generated notebooks, authenticated bounded evidence, figures,
and report source are available at
\url{https://github.com/OussemaHarrabi/SpectraQuant}.

\balance
\begin{thebibliography}{9}

\bibitem{gptq}
E. Frantar, S. Ashkboos, T. Hoefler, and D. Alistarh.
\newblock GPTQ: Accurate Post-Training Quantization for Generative Pre-trained Transformers.
\newblock \emph{arXiv:2210.17323}, 2022.

\bibitem{awq}
J. Lin et al.
\newblock AWQ: Activation-aware Weight Quantization for LLM Compression and Acceleration.
\newblock \emph{arXiv:2306.00978}, 2023.

\bibitem{smoothquant}
G. Xiao, J. Lin, M. Seznec, H. Wu, J. Demouth, and S. Han.
\newblock SmoothQuant: Accurate and Efficient Post-Training Quantization for Large Language Models.
\newblock \emph{arXiv:2211.10438}, 2022.

\bibitem{hawq}
Z. Dong et al.
\newblock HAWQ-V2: Hessian Aware Trace-Weighted Quantization of Neural Networks.
\newblock \emph{NeurIPS}, 2020.

\bibitem{loftq}
Y. Li et al.
\newblock LoftQ: LoRA-Fine-Tuning-Aware Quantization for Large Language Models.
\newblock \emph{arXiv:2310.08659}, 2023.

\bibitem{lrqat}
Y. Bondarenko et al.
\newblock LR-QAT: Low-Rank Quantization-Aware Training for LLMs.
\newblock \emph{arXiv:2406.06385}, 2024.

\bibitem{pythia}
S. Biderman et al.
\newblock Pythia: A Suite for Analyzing Large Language Models Across Training and Scaling.
\newblock \emph{ICML}, 2023.

\bibitem{c4}
C. Raffel et al.
\newblock Exploring the Limits of Transfer Learning with a Unified Text-to-Text Transformer.
\newblock \emph{JMLR}, 2020.

\bibitem{pg19}
J. W. Rae et al.
\newblock Compressive Transformers for Long-Range Sequence Modelling.
\newblock \emph{arXiv:1911.05507}, 2019.

\end{thebibliography}

\end{document}
